\documentclass[11pt,a4paper]{article}
\usepackage[utf8]{inputenc}
\usepackage[T1]{fontenc}
\usepackage{amsmath,amssymb}
\usepackage[hidelinks]{hyperref}
\usepackage{booktabs}
\usepackage[margin=1in]{geometry}

\title{RMD-Q Whitepaper\\--- One Assumption for Post-Quantum Keys and Signatures\\on Ultra-Constrained Devices}
\author{nul0 (Future Proofs Tech) \\ \texttt{futureproofs@proton.me}}
\date{September 2026 \\ v0.1 --- vision and design (companion: yellowpaper with the formal specification) \\
\small \textcopyright\ 2026 nul0 (Future Proofs Tech). Licensed under GPLv3 (see LICENSE).}

\begin{document}
\maketitle

\begin{abstract}
Post-quantum standards secure the general case. Ultra-constrained devices
--- sensors, meters, wearables, field nodes --- need cryptography that fits
in kilobytes, sips energy, and still resists quantum attackers. RMD-Q is a
research program toward that goal: a \emph{new} hardness assumption coupling
a lattice problem with a multivariate-quadratic problem over one secret,
yielding \emph{both} key exchange and signatures from \emph{one} shared
$\approx$3.9\,KB codebase. This whitepaper explains the vision, the design,
what the measurements say, and -- just as carefully -- what is not yet
proven. The formal specification lives in the companion yellowpaper.
\end{abstract}

\tableofcontents
\newpage

\section{Why: the gap after standardization}

NIST's post-quantum standards (ML-KEM, ML-DSA, SLH-DSA) are the right
answer for servers, phones, and browsers. But a battery-powered sensor
cannot comfortably carry two separate post-quantum stacks (one for key
exchange, one for signatures), each with its own math, code, and audit
surface. And the community's security eggs sit in few baskets
(module lattices plus hash-based backups) -- assumption diversity, which
NIST itself pursues in its additional-signatures process, is valuable.

RMD-Q answers both concerns at once: \textbf{one new assumption, two
primitives, one codebase.}

\section{The idea in one page}

A secret $s$ is sparse and short. Two public structures bind it:
\begin{enumerate}
\item a noisy linear relation $b = As+e$ (a lattice problem: recover $s$
  from noisy equations), and
\item a quadratic system $P(s)=0$ (a multivariate problem: find a root).
\end{enumerate}
A lattice-only attacker finds candidates that fail the quadratic check;
an MQ-only attacker finds roots that violate the noisy equations. Breaking
the scheme needs \emph{both} solvers \emph{combined} -- that jointness is
the assumption. From it we derive key encapsulation (Fujisaki-Okamoto
transform, the standard recipe also behind ML-KEM) and signatures
(Fiat-Shamir with a dual opening that proves both relations at once).

Key generation plants the quadratic system around the sampled secret, so it
always terminates (rejection sampling would never finish at scale) -- with
measured statistical uniformity and a documented leakage analysis.

\section{What the measurements say}

On commodity hardware: full key exchange in $\approx$0.25\,ms; public keys
$\approx$0.8\,KB, ciphertexts down to 800\,B compressed. Decrypt-failure
analysis gives rates below $2^{-235}$ where $2^{-138}$ is the bar, backed
by 200\,000 failure-free trials. The estimator models sit a few bits under
ML-KEM at equal dimensions (stated openly, with tolerances). Every attack
tried at small scale either broke the toy (proving the methods bite) or hit
a documented wall -- including a 114\,GB out-of-memory kill on a 16-core
box. Full tables: \texttt{docs/results/}.

\section{What is not claimed}

No NIST security category. No deployment recommendation. No proof sketches
yet (they are explicitly queued). The signature scheme is toy-grade. Side
channels beyond measured timing shape are unstudied. Anyone telling you
otherwise has not read Section~6 of the yellowpaper, which lists every
non-claim in one place.

\section{Roadmap}

Short term: signature production parameters and proofs; compressed wire
formats in C; embedded deployment diet (caller-provided buffers).
Medium term: joint-attack analysis with stronger engines; external review;
fpylll full-attack cross-checks. Long term: the assumption earns trust the
only way assumptions ever do -- by surviving years of public attack. We
publish everything needed to attempt it.

\section{Call to action}

Break it. The challenge instances, reference code, attack scripts, and the
exact list of open problems are in the repository. If the assumption falls,
we want to know first -- and the kill-phase log shows we report breaks,
including our own design bugs, as prominently as wins.
\end{document}
